\titlespacing*{\subsection}{0pt}{1.0em}{0.3em}
\titlespacing*{\paragraph}{0pt}{1.0em}{0.3em}

최근 세계 금융 시장에서 국가 간 경제적 상호 영향이 커지면서 환율 변동의 복잡성도 함께 높아지고 있다~\cite{MirandaAgrippino2020}. 환율은 개별 국가의 경제 지표만으로 설명하기 어렵고, 인접 국가 및 주요 교역국 간의 유기적인 관계 속에서 서로 영향을 주고받으며 복합적인 변동 패턴을 보인다. 그러나 전통적인 시계열 모델들은 이러한 국가 간 상호의존성을 충분히 반영하지 못하고, 개별 시계열의 과거 패턴 학습에 의존하는 한계를 보여 왔다~\cite{Fischer2018, Zhou2021Informer}. 그 결과 국가 간 금융 전이효과를 충분히 포착하기 어렵고, 장기 추세의 변화를 해석하는 데에도 제약이 남아 다양한 연구가 이루어지고 있다~\cite{Diebold2012}.

따라서 본 연구는 환율이 국가 간 상호의존적 관계로 연결되어 있다는 점에 주목하고, 이 관계를 시계열 데이터로부터 직접 학습하는 B-MTGNN(Bayesian Multivariate Time-series Graph Neural Network) 기반 예측 모형을 제안한다~\cite{Wu2020MTGNN}. 명시적 연결 구조가 없는 외환 시장의 특성을 고려하여 적응형 그래프 학습 모듈로 미국, 한국, 일본 3개국 사이의 시공간적 의존 구조를 모델링하고, 금리·물가 등 거시경제 지표를 함께 결합하여 시계열 패턴과 경제적 맥락을 동시에 반영하도록 설계하였다. 또한 수치 예측에 그치지 않고, 소형 언어 모델(Small Language Model, SLM) 기반의 멀티 에이전트 시스템(Multi-Agent System, MAS)을 결합하여 예측 결과 해석 절차를 함께 구성하였다~\cite{Yu2024FinAgent, Wu2023AutoGen}. 이를 통해 정량적 예측과 정성적 추론을 하나로 결합하고자 한다.

논문의 구성은 다음과 같다. 2장에서는 환율 예측 및 거시경제 요인에 관한 선행 연구를 고찰하고, 3장에서는 제안하는 B-MTGNN 및 멀티 에이전트 시스템의 상세 구조를 설명한다. 4장에서는 실험 과정과 결과를 통해 모델의 유효성을 검증하며, 마지막으로 5장과 6장에서는 결론 및 향후 연구 방향에 대해 논한다.
